\begin{frame}[t]{Entrada, processamento e saída}
\begin{center}
\begin{tikzpicture}[node distance=0.45cm, every node/.style={font=\small},box/.style={draw=darkgreen,fill=darkgreen!10,rounded corners,text width=2.5cm,minimum height=1.25cm,align=center}]
\node[box] (e) {Entrada\\Temperatura: 32};
\node[box,right=of e] (p) {Processamento\\Comparar com 30};
\node[box,right=of p] (s) {Saída\\Sinalizador ligado};
\draw[-{Latex},thick] (e)--(p);\draw[-{Latex},thick] (p)--(s);
\end{tikzpicture}
\end{center}
\begin{itemize}
\item No computador, os dados podem vir do teclado ou de valores de teste.
\item No microcontrolador, podem vir de sensores e botões.
\item A mesma regra pode ser exercitada sem conectar o circuito.
\end{itemize}
\end{frame}
